% ============================================================
% Preamble for LuaLaTeX + Beamer + Japanese
% Japanese: Gothic
% English: Beamer / metropolis default
% ============================================================

% --------------------
% Theme
% --------------------
\usetheme[progressbar=frametitle]{metropolis}

\setbeamertemplate{headline}[miniframes]
\setbeamertemplate{theorems}[numbered]

% --------------------
% Colors
% --------------------
\setbeamercolor{background canvas}{bg=white}
\setbeamercolor{frametitle}{fg=white, bg=teal!90!purple}
\setbeamercolor{section in head/foot}{fg=pink!70!red, bg=gray!30!black}
\setbeamercolor{alerted text}{fg=purple!70!orange}

% --------------------
% Fonts
% --------------------
% Beamer / metropolis の英語フォント設定を尊重する
\usefonttheme{professionalfonts}

% --------------------
% Japanese / English fonts for LuaLaTeX
% --------------------
\usepackage{luatexja}
\usepackage[haranoaji,deluxe]{luatexja-preset}
\usepackage{fontspec}

% English: Beamer default-like sans-serif
\setmainfont{Latin Modern Sans}
\setsansfont{Latin Modern Sans}
\setmonofont{Latin Modern Mono}

% Japanese: Gothic, like the attached slide
\renewcommand{\kanjifamilydefault}{\gtdefault}

% Make Beamer text use sans-serif
\renewcommand{\familydefault}{\sfdefault}


% --------------------
% Math packages
% --------------------
\usepackage{amsmath}
\usepackage{amssymb}
\usepackage{mathtools}
\usepackage{bm}
\usepackage{bbm}

% --------------------
% Figures / tables
% --------------------
\usepackage{graphicx}
\usepackage{booktabs}
\usepackage{multirow}

% --------------------
% Boxes / icons
% --------------------
\usepackage{tcolorbox}
\usepackage{fontawesome5}

% --------------------
% Hyperlinks
% --------------------
\hypersetup{
  unicode=true,
  colorlinks=true,
  citecolor=blue,
  linkcolor=purple!70!orange,
  urlcolor=purple!70!orange
}

% --------------------
% Theorem environments
% --------------------
\newtheorem{thm}{Theorem}[section]
\newtheorem{proposition}[thm]{Proposition}

\theoremstyle{example}
\newtheorem{exam}[thm]{Example}
\newtheorem{remark}[thm]{Remark}
\newtheorem{question}[thm]{Question}
\newtheorem{prob}[thm]{Problem}

% --------------------
% Operators and notation
% --------------------
\DeclareMathOperator{\Avar}{Avar}
\DeclareMathOperator{\varop}{var}
\DeclareMathOperator{\Var}{Var}
\DeclareMathOperator{\cov}{cov}
\DeclareMathOperator{\Cov}{Cov}
\DeclareMathOperator{\rank}{rank}
\DeclareMathOperator{\diag}{diag}
\DeclareMathOperator{\MSE}{MSE}
\DeclareMathOperator{\IMSE}{IMSE}
\DeclareMathOperator{\bias}{bias}

\DeclareMathOperator*{\argmax}{arg\,max}
\DeclareMathOperator*{\argmin}{arg\,min}

\newcommand{\E}{\mathbb{E}}
\newcommand{\R}{\mathbb{R}}
\newcommand{\Hbb}{\mathbb{H}}
\newcommand{\Pbb}{\mathbb{P}}
\newcommand{\Sbb}{\mathbb{S}}

\newcommand{\parrow}{\xrightarrow{p}}
\newcommand{\darrow}{\xrightarrow{d}}
\newcommand{\asarrow}{\xrightarrow{a.s.}}

\newcommand{\plim}{\operatorname*{plim}}
\newcommand{\iid}{\text{i.i.d.}}
\newcommand{\Normal}{\text{Normal}}
\newcommand{\Unif}{\text{Uniform}}
\newcommand{\CI}{\text{CI}}

\newcommand{\forword}{\text{ for }}
\newcommand{\vs}{\text{ vs }}
\newcommand{\where}{\text{ where }}

\newcommand{\what}[1]{\widehat{#1}}
\newcommand{\rnabla}{\rotatebox[origin=c]{180}{$\nabla$}}